\section{Introduction}\label{sec:introduction}

The minimal model program seeks a birational model whose canonical
class, or more generally whose adjoint $K_X+\Delta$, is nef. When no
such model is compatible with pseudo-effectivity, its expected outcome
is a Mori fibre space: a fibration along which the adjoint is negative.
Each birational step contracts an extremal ray or replaces a small
contraction by its flip. Constructing the individual steps and choosing
a sequence that reaches an endpoint are separate problems.

For compact Kähler fourfolds, we prove that a suitable choice of ordinary
negative steps gives a finite program. It starts on the given rational
klt pair in the global Weil-divisor $\Q$-factorial category and reaches
the endpoint dictated by pseudo-effectivity. The absence of a global
ample polarization makes both the contractions and the finiteness
argument analytic. A second issue is categorical: a divisor or reflexive
sheaf defined on the entire compact space carries information different
from a divisor germ on an arbitrary analytic neighbourhood. The proof
keeps this distinction, and keeps the actual canonical rational line
bundle, throughout the program.

\subsection{Category and the finite-program theorem}

All varieties are reduced irreducible second-countable complex analytic
spaces. A Kähler form on a normal space has smooth strictly
plurisubharmonic local potentials in local embeddings. Projectivity of a
morphism means the existence of a relatively ample holomorphic line
bundle.

\begin{definition}\label{def:global-factoriality}\label{def:strong-factoriality}
A normal compact space $X$ is \emph{globally Weil $\Q$-factorial} if
every prime Weil divisor defined on the whole of $X$ is $\Q$-Cartier and
some positive reflexive power of its canonical sheaf $\omega_X$ is a line
bundle. It is \emph{globally strongly $\Q$-factorial} if every coherent
rank-one reflexive sheaf $\cF$ on $X$ has an invertible positive reflexive
power $\cF^{[m]}=(\cF^{\otimes m})^{**}$.
\end{definition}

The Weil convention is that of \cite[Definition~2.1]{DHY2023};
the strong convention is used in \cite[Section~2]{HX2026}.
The strong condition implies the Weil
condition. Neither definition imposes $\Q$-factoriality on every analytic
local ring. The distinction matters even for projective varieties; an
example is given in Appendix~\ref{sec:local-convention}.

A rational line bundle is an element of $\Pic(X)\otimes_{\Z}\Q$.
An equality of rational line bundles means an isomorphism of holomorphic
line bundles after a common positive integral multiple.
We write $K_X$ additively for the canonical sheaf. When a canonical Weil
divisor is given, we keep that representative and transport it through
the program. We will also prove the intrinsic variant in which only the
canonical rational line bundle is specified. In that variant, identities
of adjoints mean isomorphisms of holomorphic line bundles after a common
positive multiple. Relative canonical divisors are defined using
compatible local canonical generators. In particular, a numerical
identity alone never specifies such a divisor.

For an effective rational boundary $\Delta$, put $D=K_X+\Delta$.
We use \emph{log discrepancies} throughout:
\[
 a(E;X,\Delta)
 =1+\operatorname{coeff}_E\bigl(K_W-p^*(K_X+\Delta)\bigr)
\]
on a smooth model $p:W\to X$ carrying $E$. A pair is klt when all these
numbers are positive. It is terminal when it is klt and they are greater
than one for every exceptional divisor over $X$. The same convention
will apply to the auxiliary real and generalized pairs.

Let $\BC(X)$ denote real Bott--Chern cohomology of closed $(1,1)$-forms
with local potentials, and let $\NA(X)$ be the closed cone of positive
currents of bidimension $(1,1)$ modulo their pairings with these classes.
A class is \emph{nef} if it belongs to the closure of the Kähler cone,
and \emph{pseudo-effective} if it contains a closed positive
$(1,1)$-current with local potentials. These terms for a rational line
bundle refer to its first Chern class. A class is \emph{big} if it
contains a Kähler current, meaning a closed current that dominates
a positive multiple of a Kähler form.
For a projective contraction
$f:X\to Z$, write $\rho(X/Z)$ for the dimension of the space of global
real line-bundle classes modulo zero degree on every contracted curve,
and put
\[
 \rho_{\mathrm{BC}}(X/Z)
 =\dim_{\R}\bigl(\BC(X)/f^*\BC(Z)\bigr).
\]

An \emph{ordinary $D$-negative step} is a projective divisorial
contraction, or the flip of a projective small contraction, of a
$D$-negative extremal ray of $\NA(X)$. The contraction has connected
fibres and a normal compact Kähler target; its relative rank is one and
$-D$ is relatively ample. For a flip the transformed adjoint is ample
over the same target. The boundary is pushed forward in a divisorial
step and strictly transformed in a flip. A \emph{Mori fibre space} has
the same contraction properties and a base of strictly smaller
dimension.

\begin{theorem}[Finite ordinary programs]\label{thm:finite}
Let $X$ be a normal globally Weil $\Q$-factorial compact Kähler
fourfold, with a canonical Weil divisor $K_X$, and let $\Delta\geq0$
be a rational divisor such that $(X,\Delta)$ is klt. Then there is a
finite sequence of ordinary negative steps
\[
 (X,\Delta)=(X_0,\Delta_0)\dashrightarrow\cdots
 \dashrightarrow(X_n,\Delta_n).
\]
Every $X_i$ is normal, compact Kähler and globally Weil
$\Q$-factorial; every $(X_i,\Delta_i)$ is klt, and
$D_i=K_{X_i}+\Delta_i$ is $\Q$-Cartier. The canonical divisors are
the transforms of the specified datum. The following alternatives hold.
\begin{enumerate}[label=\textup{(\roman*)}]
\item If $D=K_X+\Delta$ is pseudo-effective, then $D_n$ is nef.
The composite $\phi:X\dashrightarrow X_n$ extracts no prime divisor,
and
\[
 a(E;X,\Delta)\leq a(E;X_n,\Delta_n)
\]
for every prime divisor over the two models, with strict inequality
for every prime on $X$ contracted by $\phi$.
\item If $D$ is not pseudo-effective, there is a projective surjective
morphism $g:X_n\to S$ with connected fibres, where $S$ is normal
compact Kähler, $\dim S<4$, $\rho(X_n/S)=1$, and $-D_n$ is
$g$-ample.
\end{enumerate}
For each negative contraction, including $g$, the relative
Bott--Chern dimension is one. If $X$ is globally strongly
$\Q$-factorial, so is every $X_i$.

The same conclusions hold in the intrinsic formulation with the
canonical rational line bundle in place of a chosen canonical Weil
divisor. In either formulation the first model is exactly $X$.
\end{theorem}

Theorem~\ref{thm:finite} gives an affirmative solution of the
finite-program form of the rational klt fourfold minimal model
conjecture in this global compact Kähler category. The quantifier is
existence of one finite ordinary program; its proof uses auxiliary
models without replacing the first model $X$. In the pseudo-effective
case, the endpoint assertion is nefness, with the discrepancy properties
displayed above.

\subsection{History and the additional argument}

The projective minimal model program supplies the geometric framework:
negative extremal contractions, flips and the comparison of discrepancies.
Kawamata--Matsuda--Matsuki developed the four-dimensional terminal
termination argument using discrepancies and the loss of algebraic
surface classes~\cite[Theorem~5-1-15 and
Lemmas~5-1-16--5-1-17]{KMM1987}; Kollár--Mori give a systematic account
of the birational tools~\cite{KM1998}. Fujino extended the
fourfold termination method to canonical pairs with rational boundary
\cite{Fujino2004,Fujino2005Addendum}. His addendum identifies a
necessary correction to the low-discrepancy count: all valuations over
certain smooth boundary surfaces must be treated together, rather than
only the divisor obtained by the first blowup.

In the compact Kähler setting, Höring--Peternell constructed minimal
models for non-uniruled threefolds~\cite{HP2016}.
Das--Hacon--Păun proved the fourfold minimal model theorem for dlt
pairs whose adjoint is $\Q$-linearly equivalent to an effective divisor
\cite[Theorem~1.1]{DHP2024}. Their work also develops analytic
finite-generation and birational tools. Fujino's relative analytic
MMP supplies ordinary cone, base-point-free and flip theorems for
projective morphisms between complex analytic spaces~\cite{Fujino2022}.
These are relative projective constructions; neither the source nor
the final model is required to be projective over a point.

Generalized pairs provide a way to carry a nef class on a higher model
through the program. Das--Hacon--Yáñez developed their compact Kähler
threefold theory and the analytic language used here~\cite{DHY2023}.
Hacon--Xie's cone and transcendental base-point-free theorems supply
contractions in the compact Kähler category
\cite[Theorems~1.3--1.4]{HX2026}. Their generalized MMP with scaling
terminates when the boundary plus nef part is big
\cite[Theorem~1.5]{HX2026}; their Theorem~1.1 gives Mori fibre
models in the non-pseudo-effective case and good models in its stated
big cases. These results furnish the steps, the auxiliary programs,
and the finiteness of marked models away from scaling parameter zero
that we use below.

Matsumura--Zhong prove termination with Kähler scaling for strongly
$\Q$-factorial compact Kähler klt pairs whose pseudo-effective adjoint
has numerical dimension zero. They also obtain minimal models after
a small crepant strong $\Q$-factorial modification in that case
\cite[Theorems~6.1 and~6.3]{MatsumuraZhong2026}.

The remaining issue for Theorem~\ref{thm:finite} is the
pseudo-effective adjoint without a bigness or effectivity assumption.
We follow the limiting-model strategy of Birkar
\cite[Theorem~1.2 and Section~4]{Birkar2009}. In that strategy,
termination for a limiting terminal pair is combined with rational
decomposition of a nef real adjoint and a sufficiently small
perturbation. Birkar already works with real boundaries in the
projective category. Here the required replacements concern compact
Kähler models, potentially transcendental scaling data, and varying
projective contraction bases.

Three points make this transfer possible. First, the trace of the one
initial Kähler form is identified as a positive current with local
potentials on every relevant model. Second, we prove ordinary terminal
fourfold termination with effective real boundary by combining the
corrected low-discrepancy argument with compact analytic cycle classes.
Third, actual line bundles numerically trivial over the relevant
birational contractions descend without an additional tensor power.
The latter fact preserves the Cartier indices of finitely many rational
nef adjoints and therefore the positive intersection gap used in the
perturbation argument. None of these auxiliary constructions changes
the initial model of the ordinary program.

\subsection{Proof strategy and organization}

Sections~\ref{sec:analytic}--\ref{sec:counts} establish the common
analytic, birational and termination tools. The proof of
Theorem~\ref{thm:finite} then has four stages.
\begin{enumerate}[label=\textup{\arabic*.}]
\item Fix a Kähler form on the original space and run ordinary steps
with its transformed generalized nef data as a scaling direction.
If there were infinitely many steps, marked-model finiteness would
force their scaling parameters to tend to zero.
\item Extract the stabilized finite set of low-discrepancy places.
On a fixed auxiliary model the boundaries decrease to an effective
ordinary terminal boundary.
\item Run a finite program for that limiting terminal pair and
express its nef adjoint as a positive combination of rational nef
adjoints. Their fixed Cartier indices give a uniform positive lower
bound for nonzero degrees on curves.
\item A sufficiently small generalized perturbation has a finite
program on which all these rational adjoints are crepant. At its
endpoint, descent to an earlier flip base makes the originally
negative adjoint trivial on its contracted ray, a contradiction.
\end{enumerate}
Section~\ref{sec:endpoints} identifies the endpoint and concludes the
finite-program proof. All extractions occur on auxiliary models;
the ordinary sequence itself begins on $X$.

Section~\ref{sec:analytic} constructs the trace and comparison tools,
and Section~\ref{sec:absolute} supplies ordinary steps and exact bundle
descent. Section~\ref{sec:counts} proves the terminal theorem before
it is used. Sections~\ref{sec:scaling} and~\ref{sec:perturbation}
carry out the four stages above, and Section~\ref{sec:endpoints}
identifies the endpoints. Appendix~\ref{sec:local-convention} gives
a smooth projective example showing why analytic local
$\Q$-factoriality cannot replace the global convention in the theorem.
Figure~\ref{fig:dependencies} records the dependencies of this argument.

\begin{figure}[htbp]
\centering
\begin{tikzpicture}[
 box/.style={draw=linkblue!65,rounded corners=2pt,align=center,
 fill=linkblue!3,text width=5.1cm,inner sep=6pt,font=\small},
 edge/.style={-{Stealth[length=2mm]},draw=linkblue!80,thick}]
\node[box] (steps) at (0,0)
 {Positive traces and ordinary steps\\Exact line-bundle descent};
\node[box] (terminal) at (6.2,0)
 {Discrepancies and analytic cycle ranks\\Terminal fourfold termination\\with real boundary};
\node[box] (scaling) at (0,-2)
 {Ordinary scaling from $X$\\Vanishing thresholds\\Stabilized auxiliary extraction};
\node[box] (limit) at (6.2,-2)
 {The limiting terminal pair\\A finite small program\\to a nef ordinary adjoint};
\node[box] (gap) at (0,-4.1)
 {Rational nef decomposition\\Fixed Cartier indices\\and a positive intersection gap};
\node[box] (perturb) at (6.2,-4.1)
 {A nearby generalized program\\Crepant rational adjoints\\Descent to the original flip base};
\node[box,text width=11.3cm] (finite) at (3.1,-6.1)
 {The original negative ray becomes trivial: contradiction\\
 A finite ordinary program and its nef or Mori endpoint
 (Theorem~\ref{thm:finite})};
\draw[edge] (steps)--(scaling);
\draw[edge] (terminal)--(limit);
\draw[edge] (scaling)--(limit);
\draw[edge] (limit.south west)--(gap.north east);
\draw[edge] (gap)--(perturb);
\draw[edge] (limit)--(perturb);
\draw[edge] (perturb)--(finite);
\end{tikzpicture}
\caption{The finite-program argument. The ordinary sequence starts on
$X$. Extraction and the limiting and nearby programs are auxiliary
comparisons. Exact bundle descent preserves the indices used in the
positive intersection gap; the last descent occurs on a common
resolution, as in Figure~\ref{fig:perturbation-descent}.}
\label{fig:dependencies}
\end{figure}

